\documentclass[../../main.tex]{subfiles}
\begin{document}

\begin{center}
  {\fontsize{48pt}{58pt}\selectfont \textbf{2000}\normalsize\textbf{年度}}
\end{center}

\begin{table}[htb]
  \centering
  \caption{出題テーマと難易度}
  \begin{tabular}{|c|p{9cm}|c|}
    \hline
    \textbf{問題番号} & \textbf{テーマ} & \textbf{難易度} \\
    \hline
    第1問 & 複素数の比の虚部と三角形の内部 & \\
    \hline
    第2問 & 指数関数の不等式と最大値の級数 & \\
    \hline
    第3問 & 格子点と一次不定方程式 & \\
    \hline
    第4問 & 直方体の斜めの切断面積 & \\
    \hline
    第5問 & 複素数の累乗の大きさ & \\
    \hline
    第6問 & 逆正接関数の法線と面積 & \\
    \hline
  \end{tabular}
\end{table}

\begin{table}[htb]
  \centering
  \caption{解答}
  \renewcommand{\arraystretch}{1.4}
  \begin{tabular}{|c|c|p{9cm}|}
    \hline
    \textbf{問題番号} & \textbf{小問} & \textbf{答え} \\
    \hline
    第1問 & (1) & 領域は解答の図を参照 \\
    \cline{2-3}
    & (2) & 証明問題 \\
    \hline
    第2問 & (1) & 証明問題 \\
    \cline{2-3}
    & (2) & $\dfrac{1}{(e-1)^2}$ \\
    \hline
    第3問 & (1) & $(-ak,k)$の一例 \\
    \cline{2-3}
    & (2)，(3) & 証明問題 \\
    \hline
    第4問 & (1) & 証明問題 \\
    \cline{2-3}
    & (2) & $S(x)=
    \begin{dcases}2\sqrt{2}x^2&(0\le x\le\dfrac{1}{2})\\ \sqrt{2}\{1-2(x-1)^2\}&(\dfrac{1}{2}\le x\le\dfrac{3}{2})\\ 2\sqrt{2}(2-x)^2&(\dfrac{3}{2}\le x\le2)
    \end{dcases}$ \\
    \hline
    第5問 & (1)，(2) & 証明問題 \\
    \hline
    第6問 & (1) & $y=-2x+2+\dfrac{\pi}{4}$ \\
    \cline{2-3}
    & (2) & $\dfrac{\pi^2}{64}+\dfrac{\pi}{4}+\log\dfrac{\sqrt{2}}{2}$ \\
    \hline
  \end{tabular}
\end{table}

\paragraph{試験形態}
6問．試験時間は要確認．

\end{document}
